% GENERATED FILE. Edit canonical lessons, then run npm run generate:books.
% canonical-source-hash: fnv1a64:bd08167e002d2bee
% canonical-lessons: KA-C390-kanunu, KA-C390-kayde, KA-C390-pariskarisu, KA-C390-adhikrta, KA-C390-hastaksepa, KA-R390-review-384-390

\chapter{The Law}
\label{ch:ka-the-law}
\hlchaptermodality{390}

\begin{chapteropening}
\textbf{By the end of this chapter:} \emph{I can say a law, an act of law, to revise, official and interference.}

Complete the last lesson of chapter 390: I can say a law, an act of law, to revise, official and interference.
\end{chapteropening}

\section[\texorpdfstring{\kn{ಕಾನೂನು} (kānūnu)}{ಕಾನೂನು (kānūnu)}]{\kn{ಕಾನೂನು} (kānūnu) — a law; the law}

\label{lesson:KA-C390-kanunu}



\begin{quote}
Before the new one: say the Kannada for to exaggerate, then the Kannada for status.
\end{quote}

\subsection*{You'll want to know: \kn{ಕಾನೂನು}}
\textbf{\kn{ಕಾನೂನು}} — \emph{kānūnu} — \textquotedblleft{}a law; the law\textquotedblright{}.

\begin{quote}
\kn{ಈ} \kn{ಕಾನೂನು} \kn{ಮುಂದಿನ} \kn{ತಿಂಗಳಿನಿಂದ} \kn{ಬರುತ್ತದೆ}.
\end{quote}

\begin{quote}
\emph{ī kānūnu mundina tiṅgaḷininda baruttade.}
\end{quote}

\begin{quote}
This law comes in from next month.
\end{quote}

\kn{ಕಾನೂನುಬದ್ಧ}, legal; \kn{ಕಾನೂನುಬಾಹಿರ}, illegal.

\begin{cousinweb}
\textbf{\kn{ಕಾನೂನು}} (\emph{kānūnu}) is a loan from Arabic \emph{qānūn}, \textquotedblleft{}a rule, a canon\textquotedblright{}, through Persian and Hindustani; Arabic took it from Greek \emph{kanōn}, \textquotedblleft{}a rule\textquotedblright{}.
\end{cousinweb}

\begin{grammarlens}[title={What you've built}]
One more word for giving reasons and opinions.
\end{grammarlens}

\subsection*{Your turn}
\begin{itemize}
\raggedright
  \item \emph{Say it:} \emph{kānūnu}
  \item \emph{Say it:} \emph{kānūnu}, once more
  \item \emph{Say it:} \emph{ī kānūnu mundina tiṅgaḷininda baruttade.} — This law comes in from next month.
  \item \emph{Recall:} say the Kannada for to exaggerate, then the Kannada for status, then say \emph{kānūnu} again
\end{itemize}

\begin{tcolorbox}[breakable,colback=teal!4,colframe=teal!35!black,title={Before you move on}]
What does \kn{ಕಾನೂನು} mean? (\textquotedblleft{}A law; the law\textquotedblright{}.) Say it once more.
\end{tcolorbox}
\section[\texorpdfstring{\kn{ಕಾಯ್ದೆ} (kāyde)}{ಕಾಯ್ದೆ (kāyde)}]{\kn{ಕಾಯ್ದೆ} (kāyde) — an act of law}

\label{lesson:KA-C390-kayde}



\begin{quote}
Before the new one: say the Kannada for status, then the Kannada for a law.
\end{quote}

\subsection*{You'll want to know: \kn{ಕಾಯ್ದೆ}}
\textbf{\kn{ಕಾಯ್ದೆ}} — \emph{kāyde} — \textquotedblleft{}an act of law\textquotedblright{}.

\begin{quote}
\kn{ಸರ್ಕಾರ} \kn{ಹೊಸ} \kn{ಕಾಯ್ದೆ} \kn{ತಂದಿತು}.
\end{quote}

\begin{quote}
\emph{sarkāra hosa kāyde tanditu.}
\end{quote}

\begin{quote}
The government brought in a new act.
\end{quote}

A law passed under a name: \kn{ಕಾಯ್ದೆ}. The law in general: \kn{ಕಾನೂನು}.

\begin{cousinweb}
\textbf{\kn{ಕಾಯ್ದೆ}} (\emph{kāyde}) is a loan from Arabic \emph{qāʿida}, \textquotedblleft{}a rule, a base\textquotedblright{}, which reached Kannada through Hindustani.
\end{cousinweb}

\begin{grammarlens}[title={What you've built}]
One more word for giving reasons and opinions.
\end{grammarlens}

\subsection*{Your turn}
\begin{itemize}
\raggedright
  \item \emph{Say it:} \emph{kāyde}
  \item \emph{Say it:} \emph{kāyde}, once more
  \item \emph{Say it:} \emph{sarkāra hosa kāyde tanditu.} — The government brought in a new act.
  \item \emph{Recall:} say the Kannada for status, then the Kannada for a law, then say \emph{kāyde} again
\end{itemize}

\begin{tcolorbox}[breakable,colback=teal!4,colframe=teal!35!black,title={Before you move on}]
What does \kn{ಕಾಯ್ದೆ} mean? (\textquotedblleft{}An act of law\textquotedblright{}.) Say it once more.
\end{tcolorbox}
\section[\texorpdfstring{\kn{ಪರಿಷ್ಕರಿಸು} (pariṣkarisu)}{ಪರಿಷ್ಕರಿಸು (pariṣkarisu)}]{\kn{ಪರಿಷ್ಕರಿಸು} (pariṣkarisu) — to revise}

\label{lesson:KA-C390-pariskarisu}



\begin{quote}
Before the new one: say the Kannada for a law, then the Kannada for an act of law.
\end{quote}

\subsection*{You'll want to know: \kn{ಪರಿಷ್ಕರಿಸು}}
\textbf{\kn{ಪರಿಷ್ಕರಿಸು}} — \emph{pariṣkarisu} — \textquotedblleft{}to revise\textquotedblright{}.

\begin{quote}
\kn{ಸರ್ಕಾರ} \kn{ಹಳೆಯ} \kn{ನಿಯಮಗಳನ್ನು} \kn{ಪರಿಷ್ಕರಿಸಿತು}.
\end{quote}

\begin{quote}
\emph{sarkāra haḷeya niyamagaḷannu pariṣkarisitu.}
\end{quote}

\begin{quote}
The government revised the old rules.
\end{quote}

Past (I …): \kn{ಪರಿಷ್ಕರಿಸಿದೆ}, \emph{pariṣkariside}.

\kn{ಪರಿಷ್ಕೃತ}, revised.

\begin{cousinweb}
\textbf{\kn{ಪರಿಷ್ಕರಿಸು}} (\emph{pariṣkarisu}) is built on Sanskrit \emph{pariṣkāra}, \textquotedblleft{}adorning, preparing, refining\textquotedblright{} (CDIAL 7894), with the Kannada verb ending -\emph{isu}.
\end{cousinweb}

\begin{grammarlens}[title={What you've built}]
One more word for giving reasons and opinions.
\end{grammarlens}

\subsection*{Your turn}
\begin{itemize}
\raggedright
  \item \emph{Say it:} \emph{pariṣkarisu}
  \item \emph{Say it:} \emph{pariṣkarisu}, once more
  \item \emph{Say it:} \emph{sarkāra haḷeya niyamagaḷannu pariṣkarisitu.} — The government revised the old rules.
  \item \emph{Recall:} say the Kannada for a law, then the Kannada for an act of law, then say \emph{pariṣkarisu} again
\end{itemize}

\begin{tcolorbox}[breakable,colback=teal!4,colframe=teal!35!black,title={Before you move on}]
What does \kn{ಪರಿಷ್ಕರಿಸು} mean? (\textquotedblleft{}To revise\textquotedblright{}.) Say it once more.
\end{tcolorbox}
\section[\texorpdfstring{\kn{ಅಧಿಕೃತ} (adhikṛta)}{ಅಧಿಕೃತ (adhikṛta)}]{\kn{ಅಧಿಕೃತ} (adhikṛta) — official}

\label{lesson:KA-C390-adhikrta}



\begin{quote}
Before the new one: say the Kannada for an act of law, then the Kannada for to revise.
\end{quote}

\subsection*{You'll want to know: \kn{ಅಧಿಕೃತ}}
\textbf{\kn{ಅಧಿಕೃತ}} — \emph{adhikṛta} — \textquotedblleft{}official\textquotedblright{}.

\begin{quote}
\kn{ಇದು} \kn{ಅಧಿಕೃತ} \kn{ಸುದ್ದಿ} \kn{ಅಲ್ಲ}.
\end{quote}

\begin{quote}
\emph{idu adhikṛta suddi alla.}
\end{quote}

\begin{quote}
This is not official news.
\end{quote}

\kn{ಅಧಿಕೃತವಾಗಿ}, officially.

\begin{cousinweb}
\textbf{\kn{ಅಧಿಕೃತ}} (\emph{adhikṛta}) is borrowed from Sanskrit \emph{adhikṛta}, \textquotedblleft{}placed at the head of, appointed\textquotedblright{} (Monier-Williams).
\end{cousinweb}

\begin{grammarlens}[title={What you've built}]
One more word for giving reasons and opinions.
\end{grammarlens}

\subsection*{Your turn}
\begin{itemize}
\raggedright
  \item \emph{Say it:} \emph{adhikṛta}
  \item \emph{Say it:} \emph{adhikṛta}, once more
  \item \emph{Say it:} \emph{idu adhikṛta suddi alla.} — This is not official news.
  \item \emph{Recall:} say the Kannada for an act of law, then the Kannada for to revise, then say \emph{adhikṛta} again
\end{itemize}

\begin{tcolorbox}[breakable,colback=teal!4,colframe=teal!35!black,title={Before you move on}]
What does \kn{ಅಧಿಕೃತ} mean? (\textquotedblleft{}Official\textquotedblright{}.) Say it once more.
\end{tcolorbox}
\section[\texorpdfstring{\kn{ಹಸ್ತಕ್ಷೇಪ} (hastakṣēpa)}{ಹಸ್ತಕ್ಷೇಪ (hastakṣēpa)}]{\kn{ಹಸ್ತಕ್ಷೇಪ} (hastakṣēpa) — interference}

\label{lesson:KA-C390-hastaksepa}



\begin{quote}
Before the new one: say the Kannada for to revise, then the Kannada for official.
\end{quote}

\subsection*{You'll want to know: \kn{ಹಸ್ತಕ್ಷೇಪ}}
\textbf{\kn{ಹಸ್ತಕ್ಷೇಪ}} — \emph{hastakṣēpa} — \textquotedblleft{}interference\textquotedblright{}.

\begin{quote}
\kn{ನಮ್ಮ} \kn{ಮನೆಯ} \kn{ವಿಷಯದಲ್ಲಿ} \kn{ಬೇರೆಯವರ} \kn{ಹಸ್ತಕ್ಷೇಪ} \kn{ಬೇಡ}.
\end{quote}

\begin{quote}
\emph{namma maneya viṣayadalli bēreyavara hastakṣēpa bēḍa.}
\end{quote}

\begin{quote}
We want no outside interference in our family's matters.
\end{quote}

Word by word it is throwing in a hand: \kn{ಹಸ್ತ} + \kn{ಕ್ಷೇಪ}.

\begin{cousinweb}
\textbf{\kn{ಹಸ್ತಕ್ಷೇಪ}} (\emph{hastakṣēpa}) is built from Sanskrit \emph{hasta}, \textquotedblleft{}a hand\textquotedblright{} (CDIAL 14024), and \emph{kṣepa}, \textquotedblleft{}a throw, a cast\textquotedblright{} (Monier-Williams).
\end{cousinweb}

\begin{grammarlens}[title={What you've built}]
One more word for giving reasons and opinions.
\end{grammarlens}

\subsection*{Your turn}
\begin{itemize}
\raggedright
  \item \emph{Say it:} \emph{hastakṣēpa}
  \item \emph{Say it:} \emph{hastakṣēpa}, once more
  \item \emph{Say it:} \emph{namma maneya viṣayadalli bēreyavara hastakṣēpa bēḍa.} — We want no outside interference in our family's matters.
  \item \emph{Recall:} say the Kannada for to revise, then the Kannada for official, then say \emph{hastakṣēpa} again
\end{itemize}

\begin{tcolorbox}[breakable,colback=teal!4,colframe=teal!35!black,title={Before you move on}]
What does \kn{ಹಸ್ತಕ್ಷೇಪ} mean? (\textquotedblleft{}Interference\textquotedblright{}.) Say it once more.
\end{tcolorbox}
\section[(dialogue)]{Looking back: from Grasping an Idea to The Law}

\label{lesson:KA-R390-review-384-390}



\begin{quote}
Say the Kannada for to grasp and an interpretation.
\end{quote}

\subsection*{Your turn}
Say these aloud:

\begin{itemize}
\raggedright
  \item to grasp, an interpretation, a notion, an aspect, a feature
  \item influence, to affect badly, to come about, a tendency, a source
  \item an explanation, a lame excuse, accidental, to portray, blame
  \item a challenge, to set aside, balance, a priority, a value
  \item ethical, a sacrifice, a conscience, well-being, to condemn
  \item to commend, to felicitate, prestige, to exaggerate, status
  \item a law, an act of law, to revise, official, interference
\end{itemize}

\begin{tcolorbox}[breakable,colback=teal!4,colframe=teal!35!black,title={Before you move on}]
To grasp? (\textbf{\kn{ಗ್ರಹಿಸು}}, \emph{grahisu}.) An interpretation? (\textbf{\kn{ವ್ಯಾಖ್ಯಾನ}}, \emph{vyākhyāna}.) A notion? (\textbf{\kn{ಕಲ್ಪನೆ}}, \emph{kalpane}.) An aspect? (\textbf{\kn{ಅಂಶ}}, \emph{aṃśa}.) A feature? (\textbf{\kn{ಲಕ್ಷಣ}}, \emph{lakṣaṇa}.) Influence? (\textbf{\kn{ಪ್ರಭಾವ}}, \emph{prabhāva}.) To affect badly? (\textbf{\kn{ಬಾಧಿಸು}}, \emph{bādhisu}.) To come about? (\textbf{\kn{ಉಂಟಾಗು}}, \emph{uṇṭāgu}.) A tendency? (\textbf{\kn{ಪ್ರವೃತ್ತಿ}}, \emph{pravṛtti}.) A source? (\textbf{\kn{ಮೂಲ}}, \emph{mūla}.) An explanation? (\textbf{\kn{ಸಮಜಾಯಿಷಿ}}, \emph{samajāyiṣi}.) A lame excuse? (\textbf{\kn{ಸಬೂಬು}}, \emph{sabūbu}.) Accidental? (\textbf{\kn{ಆಕಸ್ಮಿಕ}}, \emph{ākasmika}.) To portray? (\textbf{\kn{ಬಿಂಬಿಸು}}, \emph{bimbisu}.) Blame? (\textbf{\kn{ಹೊಣೆ}}, \emph{hoṇe}.) A challenge? (\textbf{\kn{ಸವಾಲು}}, \emph{savālu}.) To set aside? (\textbf{\kn{ಮೀಸಲಿಡು}}, \emph{mīsaliḍu}.) Balance? (\textbf{\kn{ಸಮತೋಲನ}}, \emph{samatōlana}.) A priority? (\textbf{\kn{ಆದ್ಯತೆ}}, \emph{ādyate}.) A value? (\textbf{\kn{ಮೌಲ್ಯ}}, \emph{maulya}.) Ethical? (\textbf{\kn{ನೈತಿಕ}}, \emph{naitika}.) A sacrifice? (\textbf{\kn{ತ್ಯಾಗ}}, \emph{tyāga}.) A conscience? (\textbf{\kn{ಮನಸ್ಸಾಕ್ಷಿ}}, \emph{manassākṣi}.) Well-being? (\textbf{\kn{ಹಿತ}}, \emph{hita}.) To condemn? (\textbf{\kn{ಖಂಡಿಸು}}, \emph{khaṇḍisu}.) To commend? (\textbf{\kn{ಶ್ಲಾಘಿಸು}}, \emph{ślāghisu}.) To felicitate? (\textbf{\kn{ಪುರಸ್ಕರಿಸು}}, \emph{puraskarisu}.) Prestige? (\textbf{\kn{ಪ್ರತಿಷ್ಠೆ}}, \emph{pratiṣṭhe}.) To exaggerate? (\textbf{\kn{ಉತ್ಪ್ರೇಕ್ಷಿಸು}}, \emph{utprēkṣisu}.) Status? (\textbf{\kn{ಸ್ಥಾನಮಾನ}}, \emph{sthānamāna}.) A law? (\textbf{\kn{ಕಾನೂನು}}, \emph{kānūnu}.) An act of law? (\textbf{\kn{ಕಾಯ್ದೆ}}, \emph{kāyde}.) To revise? (\textbf{\kn{ಪರಿಷ್ಕರಿಸು}}, \emph{pariṣkarisu}.) Official? (\textbf{\kn{ಅಧಿಕೃತ}}, \emph{adhikṛta}.) Interference? (\textbf{\kn{ಹಸ್ತಕ್ಷೇಪ}}, \emph{hastakṣēpa}.)
\end{tcolorbox}
